\section{Enformer: predicción de la regulación de largo alcance y de la expresión génica a partir de la secuencia de DNA}

DeepConsensus mejora primero la pregunta de «si la secuencia que se observa es correcta»; Enformer pregunta después «qué efecto regulador producirá esta secuencia». Muchas variantes asociadas con enfermedades se encuentran en la non-coding region: no alteran directamente la codificación de proteínas, pero pueden influir en la transcription factor binding, la enhancer--promoter interaction y la gene expression. Para predecir este efecto, el modelo necesita abarcar un contexto que va mucho más allá del motif local.

\subsection{Por qué los GWAS llevan el problema hacia las regiones no codificantes}

La diapositiva del artículo presenta Enformer, y la diapositiva de GWAS muestra que, a medida que crece la cohort, el número de disease-associated loci descubiertos aumenta con rapidez. La diapositiva Coding versus non-coding expone el hecho fundamental: una gran cantidad de señales de asociación se ubica en regiones que no codifican proteínas. Estas pueden influir en los genes mediante la regulación, en lugar de alterar directamente un amino acid; por ello, los flujos de trabajo tradicionales que solo examinan las coding mutation no pueden explicar la mayoría de las asociaciones.

\lecturefigure{slide-69-gene-expression-paper.jpg}{Trabajo asignado: Effective Gene Expression Prediction from Sequence}{Diapositiva V069 recuperada de la grabación oficial; Stanford CS25 Lecture 17}

\lecturefigure{slide-70-gwas-growth.jpg}{Crecimiento de la escala de los GWAS y descubrimiento de loci asociados}{Diapositiva V070 recuperada de la grabación oficial; Stanford CS25 Lecture 17}

\lecturefigure{slide-71-coding-vs-noncoding.jpg}{Diferencias funcionales entre variantes coding y non-coding}{Diapositiva V071 recuperada de la grabación oficial; Stanford CS25 Lecture 17}

\begin{knowledgebox}{Concepto previo: GWAS, variant y causalidad}
Un \term{GWAS} compara la asociación entre las genetic variant y el fenotipo en una población; una \term{variant} es una diferencia de bases de un individuo respecto del genoma de referencia; una \term{association} indica una correlación estadística, que no equivale automáticamente a una causal variant. El desequilibrio de ligamiento, la estructura poblacional y los factores ambientales pueden hacer que varios loci cercanos resulten significativos a la vez; por eso las predicciones del modelo se usan principalmente para establecer prioridades y formular hipótesis sobre mecanismos.
\end{knowledgebox}

\subsection{Enhancer, promoter y acción a distancia}

La diapositiva Transcription regulation muestra la relación entre promoter, enhancer, transcription factor y gene expression. La diapositiva Variant disrupts binding explica además que el cambio de una sola base no codificante puede alterar la unión de un factor y, en consecuencia, modificar la expresión de un gen lejano. La convolución local puede reconocer un motif, pero difícilmente conecta de forma directa un enhancer y un promoter separados por decenas de miles de bases; esa es precisamente la motivación para usar attention en el genoma.

\lecturefigure{slide-72-transcription-regulation.jpg}{Regulación de la transcripción: promoter, enhancer y transcription factor}{Diapositiva V072 recuperada de la grabación oficial; Stanford CS25 Lecture 17}

\lecturefigure{slide-73-variant-disrupts-binding.jpg}{Cómo una variante no codificante interrumpe la unión y modifica la expresión}{Diapositiva V073 recuperada de la grabación oficial; Stanford CS25 Lecture 17}

\begin{warningbox}{Predecir la regulación no es diagnosticar una enfermedad}
El modelo puede predecir que cierto allele probablemente modifique un track de chromatin o de expression, pero el fenotipo de una enfermedad también depende del tejido, la etapa del desarrollo, el ambiente, múltiples genes y el contexto clínico. Las puntuaciones de efecto de variantes sirven para seleccionar candidatos experimentales; no deben traducirse directamente en una probabilidad individual de padecer la enfermedad.
\end{warningbox}

\subsection{Entrada, salida y arquitectura híbrida de Enformer}

La diapositiva de la tarea toma una DNA sequence larga como entrada y predice numerosos genomic tracks; la diapositiva de Enformer muestra la combinación de convolution y Transformer. La convolución comprime primero los motif locales, la attention establece después interacciones entre posiciones en un rango más amplio y las cabezas de salida predicen señales para distintos tipos celulares y mediciones experimentales. El modelo lee un contexto de aproximadamente 200,000 bases, por lo que tiene más posibilidades que los modelos de ventana corta de conectar elementos reguladores distantes.

\lecturefigure{slide-74-dna-sequence-prediction-task.jpg}{Predicción de múltiples señales genómicas funcionales a partir de la secuencia de DNA}{Diapositiva V074 recuperada de la grabación oficial; Stanford CS25 Lecture 17}

\lecturefigure{slide-75-enformer-model.jpg}{Enformer: modelo de largo alcance que combina convolución y Transformer}{Diapositiva V075 recuperada de la grabación oficial; Stanford CS25 Lecture 17}

El modelo puede abstraerse como
$$
\hat Z=f_{\theta}(X),\qquad
X\in\{A,C,G,T\}^{L},\quad
\hat Z\in\mathbb{R}^{B\times K}.
$$
donde $X$ es la secuencia de DNA de longitud $L$, $B$ es el número de bins de posición de salida, $K$ es el número de assay, cell type o genomic track distintos, y $\hat Z$ es la señal predicha. El effect de una variante puede expresarse como la diferencia entre las dos secuencias de allele:
$$
\Delta_k=f_{\theta}(X_{\text{alt}})_k-f_{\theta}(X_{\text{ref}})_k.
$$
donde $X_{\text{ref}}$ y $X_{\text{alt}}$ difieren únicamente en la variant candidata, y $\Delta_k$ es el cambio predicho en el track $k$. Esta diferencia es una counterfactual estimate dentro del modelo, no una prueba experimental de causalidad.

La arquitectura híbrida también refleja una división del trabajo por escalas. La convolución tiene un sesgo inductivo fuerte hacia los motif locales y puede reconocer con eficiencia patrones cortos de unión de factores de transcripción; el pooling amplía gradualmente el campo receptivo y reduce la longitud de la secuencia; el Transformer permite después la interacción entre posiciones lejanas ya comprimidas. Aplicar attention a resolución completa sobre 200,000 bases desde el inicio sería muy costoso; si solo se apilaran convoluciones, la señal de un enhancer lejano tendría que atravesar muchas capas para llegar al promoter. Enformer no es «la attention que sustituye a la convolución», sino una asignación de los patrones locales y de la regulación global a los componentes adecuados.

\subsection{Thousands of tracks: la salida no es una sola etiqueta}

La diapositiva Tracks muestra que el modelo debe predecir simultáneamente una gran cantidad de señales experimentales, por ejemplo la chromatin accessibility, la transcription factor binding y la gene expression en distintos tipos celulares. La salida multitarea permite que la representación compartida absorba procesos biológicos relacionados, pero también aumenta el label noise y el desequilibrio de los datos: algunos track tienen datos abundantes, mientras que ciertos tejidos y poblaciones están claramente subrepresentados.

\lecturefigure{slide-76-enformer-tracks.jpg}{Genomic tracks de múltiples tipos celulares predichos por Enformer}{Diapositiva V076 recuperada de la grabación oficial; Stanford CS25 Lecture 17}

\begin{knowledgebox}{Lectura de la figura: entender un genomic track como «función de la posición»}
El eje horizontal es la posición en el genoma y el eje vertical es la intensidad de cierta señal experimental; los distintos colores o filas representan assay, cell type o valores del modelo frente a valores reales. Primero hay que ver si las posiciones de los picos coinciden, luego si las intensidades concuerdan y, por último, verificar si un cambio en un enhancer lejano se corresponde con una modificación de la expression cerca del promoter. Mirar solo el coeficiente de correlación global puede ocultar picos raros pero decisivos.
\end{knowledgebox}

\subsection{Resultados y límites de uso}

La diapositiva de resultados compara las predicciones relacionadas con promoter y enhancer, y muestra que un contexto más largo mejora la regulatory activity prediction. La diapositiva Application takeaways resume los usos: explicar la non-coding variation, predecir la gene expression y apoyar la variant prioritization. El valor científico del modelo consiste en reducir el espacio de búsqueda experimental: identificar, entre decenas de miles de loci asociados, los candidatos que más merecen un assay.

\lecturefigure{slide-77-promoter-enhancer-results.jpg}{Mejora de Enformer en la predicción de la promoter y enhancer activity}{Diapositiva V077 recuperada de la grabación oficial; Stanford CS25 Lecture 17}

\lecturefigure{slide-78-application-takeaways.jpg}{Conclusiones de aplicación y limitaciones de Enformer}{Diapositiva V078 recuperada de la grabación oficial; Stanford CS25 Lecture 17}

\teachervoice{El ponente establece primero la necesidad con «la gran mayoría de los GWAS hits están en la non-coding region» y después explica la enhancer--promoter interaction a larga distancia. Este orden es importante: la attention no se usa porque esté de moda, sino porque el mecanismo biológico plantea un problema que excede el receptive field local.}

\begin{importantbox}{La cadena de evidencia correcta de Enformer}
Un contexto largo de DNA mejora la predicción de múltiples genomic track y de la expression; la diferencia entre el allele reference/alternate permite ordenar las variant; los experimentos verifican después la unión, la expresión y el fenotipo. El modelo vuelve más inteligente la decisión de «dónde experimentar», pero no elimina los experimentos ni convierte automáticamente una statistical association en un disease mechanism.
\end{importantbox}

Si se quiere usar Enformer para la variant prioritization, también hay que elegir los track que correspondan al tejido de la enfermedad. En una enfermedad cardiaca, examinar solo la puntuación global de una línea celular cualquiera puede producir candidatos biológicamente irrelevantes; en variantes poblacionales, además, hay que comprobar si los datos de entrenamiento representan la ancestry objetivo. Un informe más sólido enumera la puntuación del modelo, los cell type relevantes, los gene cercanos, la evidencia eQTL/epigenomic existente y la viabilidad experimental, para que los investigadores conozcan el fundamento del orden. Aquí el modelo desempeña el papel de un evidence aggregator, no de un juez final.

\subsection{Resumen de la sección}

Enformer transforma secuencias largas de DNA en genomic tracks de múltiples tipos celulares y múltiples assay; usa convoluciones para extraer motif locales y un Transformer para establecer relaciones reguladoras de largo alcance. Puede mejorar la predicción de la gene-expression y del variant-effect, y ofrece una base para ordenar los experimentos sobre variantes no codificantes; sin embargo, asociación, predicción y causalidad siguen siendo tres niveles distintos de evidencia.
